\section{Week 5: Prototyping Human-Centred AI Systems}

\subsection{Topic Summary}

\begin{keybox}[Learning Objectives]
By the end of this week, you will be able to:
\begin{itemize}
  \item \textbf{HAI2}: Describe the human-AI interaction design process and its distinction from conventional human-centred design.
  \item \textbf{HAI4}: Create and evaluate human-centred human-AI prototypes.
\end{itemize}
\end{keybox}

\begin{keybox}[Big Picture]
This is the ``doing'' week---prototyping is learning by making.

The core principle: \textbf{Maximize the rate of learning by minimizing the time to try new ideas.}

We cover:
\begin{enumerate}
  \item \textbf{What prototypes are}: Definition, purpose, and types.
  \item \textbf{Low vs.\ high fidelity}: When to use each.
  \item \textbf{Prototyping techniques}: Sketches, wireframes, paper, Figma, Wizard of Oz.
  \item \textbf{AI-specific prototyping}: How to prototype AI systems before building them.
\end{enumerate}
\end{keybox}

\begin{keybox}[What Examiners Like to Ask]
\begin{itemize}
  \item Define what a prototype is and explain its purpose.
  \item Compare low-fidelity and high-fidelity prototypes with advantages/disadvantages.
  \item Explain when you would use vertical vs.\ horizontal prototyping.
  \item Describe Wizard of Oz prototyping and give an AI example.
  \item Given an AI use case, propose an appropriate prototyping strategy.
  \item Explain why prototyping is particularly important for human-centered AI.
\end{itemize}
\end{keybox}

\subsection{Overview + Core Concepts + Extended Theory}

\subsubsection*{What are Prototypes?}

\begin{defbox}[Prototype]
One \textbf{manifestation of a design} that allows stakeholders to interact with it and explore its suitability.

\textbf{Key property}: A prototype \textbf{emphasizes one set of characteristics} and de-emphasizes others. It is intentionally \textbf{limited}.
\end{defbox}

\paragraph*{What Prototypes Can Be}
\begin{itemize}
  \item Paper-based storyboard
  \item Cardboard mockup
  \item Wireframe sketch
  \item Click-through interactive demo
  \item Partial software implementation
  \item Physical 3D-printed model
\end{itemize}

\paragraph*{Purpose of Prototypes}
\begin{enumerate}
  \item Help you \textbf{think through} possibilities in your design space.
  \item Help you \textbf{test assumptions} about users, experience, or technology viability.
  \item Help users \textbf{envision} interactions with the new system.
  \item Serve as \textbf{probes} to get feedback from users.
\end{enumerate}

\subsubsection*{Why Prototype Human-Centered AI?}

\begin{center}
\renewcommand{\arraystretch}{1.3}
\begin{tabular}{@{} l p{5cm} p{5cm} @{}}
\toprule
& \textbf{Traditional AI} & \textbf{Human-Centered AI} \\
\midrule
Focus & Model-centric & User-centric \\
Starting point & Data first & People first \\
Evaluation & Accuracy & Usability, trust, fairness \\
\bottomrule
\end{tabular}
\end{center}

\paragraph*{Reasons to Prototype HAI Systems}
\begin{itemize}
  \item \textbf{Align behavior with expectations}: Does the AI do what users expect?
  \item \textbf{Catch usability issues early}: Before expensive implementation.
  \item \textbf{Explore ethical concerns}: Trust, explainability, fairness---before deployment.
\end{itemize}

\subsubsection*{Low-Fidelity vs.\ High-Fidelity Prototypes}

\begin{center}
\renewcommand{\arraystretch}{1.3}
\begin{tabular}{@{} l p{5.5cm} p{5.5cm} @{}}
\toprule
& \textbf{Low-Fidelity} & \textbf{High-Fidelity} \\
\midrule
Cost & Lower development cost & Time/resource intensive \\
Speed & Fast to create and change & Slower to develop \\
Detail & Little detail, conceptual & Look \& feel resembles final \\
Interactivity & Limited or none & Partially/fully functional \\
Best for & Early ideation, exploring concepts & Usability testing, refinement \\
Risk & No error checking possible & May confuse stakeholders \\
Examples & Sketches, wireframes, paper & Figma, coded prototype \\
\bottomrule
\end{tabular}
\end{center}

\paragraph*{The Fidelity Spectrum}

\begin{center}
\textbf{Sketches} $\to$ \textbf{Wireframes} $\to$ \textbf{Interactive Prototype} $\to$ \textbf{Final Design}
\end{center}

\begin{itemize}
  \item \textbf{Sketches}: Low commitment, quick to throw away, good for \textit{communicating ideas}.
  \item \textbf{Wireframes}: Structure and flow, good for \textit{concept validation}.
  \item \textbf{Interactive prototypes}: Some interactivity, good for \textit{refinement and testing}.
  \item \textbf{Final design}: Fully branded, good for \textit{wide-scale evaluation}.
\end{itemize}

\begin{tipbox}[Match Fidelity to Purpose]
Don't build a high-fidelity prototype just to test an idea. Use sketches or wireframes for early exploration; save high-fidelity for later validation.
\end{tipbox}

\subsubsection*{Low-Fidelity Prototyping Techniques}

\paragraph*{1. Wireframes}
\begin{itemize}
  \item Draw overviews of interactive products to establish structure and flow.
  \item Primarily \textbf{tools for collaboration}---make better prototypes faster.
  \item Tools: \textbf{Balsamiq}, Miro, Figma (low-fi mode).
  \item Use boxes with X for images, placeholder text for content.
\end{itemize}

\paragraph*{2. Paper Prototyping}
\begin{itemize}
  \item Sketch concepts/flows on paper; test with users.
  \item Use cardboard ``phones'' or ``tablets'' to simulate device context.
  \item \textbf{Why paper?} Low stakes $\to$ users feel more comfortable critiquing.
  \item Research shows hand-drawn prototypes invite more honest feedback.
\end{itemize}

\paragraph*{3. Click-Through Prototypes (Low-Fi Version)}
\begin{itemize}
  \item Connect frames in Figma/Miro with simple click interactions.
  \item Test basic flows without full implementation.
  \item Example: Button changes screen color---tests interaction, not content.
\end{itemize}

\subsubsection*{High-Fidelity Prototyping Techniques}

\paragraph*{1. Click-Through Prototypes (High-Fi)}
\begin{itemize}
  \item Tools: \textbf{Figma}, Adobe XD, Axure.
  \item Interactive, resembles final product.
  \item Can test on mobile devices (Figma Mirror app).
  \item Start with Miro wireframes $\to$ build up in Figma.
\end{itemize}

\paragraph*{2. Vertical vs.\ Horizontal Prototyping}

\begin{center}
\renewcommand{\arraystretch}{1.3}
\begin{tabular}{@{} l p{5.5cm} p{5.5cm} @{}}
\toprule
& \textbf{Vertical} & \textbf{Horizontal} \\
\midrule
Implementation & Few functions fully implemented & All UI implemented but shallow \\
Depth & Deep (one feature works end-to-end) & Shallow (no real data processing) \\
Example & Messaging works; nothing else & All buttons visible; nothing saves \\
Best for & Testing specific feature in depth & Testing overall navigation/layout \\
\bottomrule
\end{tabular}
\end{center}

\paragraph*{3. Physical Computing}
\begin{itemize}
  \item Platforms: Arduino, Raspberry Pi, 3D printing, laser cutting.
  \item For tangible interfaces, IoT devices, wearables.
  \item Quickly create 3D form factors to test physical interactions.
\end{itemize}

\subsubsection*{Wizard of Oz Prototyping}

\begin{defbox}[Wizard of Oz]
A prototyping technique where a \textbf{human simulates} the system's behavior behind the scenes, while the user believes they are interacting with a real system.

Named after the film: A powerful wizard is actually just a person behind a curtain.
\end{defbox}

\paragraph*{Classic Example: IBM Listening Typewriter (1984)}
\begin{itemize}
  \item User speaks into microphone; text appears on screen.
  \item User believes it's speech-to-text technology.
  \item \textbf{Reality}: A human typist in another room transcribes in real-time.
  \item \textbf{Purpose}: Test the \textit{experience} before building the technology.
\end{itemize}

\paragraph*{Why Use Wizard of Oz?}
\begin{itemize}
  \item Test experiences for technology that doesn't exist yet.
  \item Test AI interactions before training/fine-tuning models.
  \item Much faster and cheaper than building real AI.
  \item Especially useful for niche domains (healthcare, hospitality, etc.).
\end{itemize}

\subsubsection*{AI-Specific Prototyping Strategies}

\begin{center}
\renewcommand{\arraystretch}{1.3}
\begin{tabular}{@{} l p{8cm} @{}}
\toprule
\textbf{AI Use Case} & \textbf{Prototype Strategy} \\
\midrule
Recommender system & Mock interface with fake suggestions (Figma/paper). \\
Chatbot / LLM & Wizard of Oz with scripted/human responses. \\
Classifier UI & Confidence bars + explanation cards. \\
Personalization & Multiple mockups showing different personalized states. \\
Agent/assistant & Wizard of Oz with human behind the scenes. \\
\bottomrule
\end{tabular}
\end{center}

\begin{exbox}[Prototyping a Chatbot]
\textbf{Scenario}: You're designing an AI tutor chatbot.

\textbf{Before building any AI}:
\begin{enumerate}
  \item Create Figma mockups of the chat interface.
  \item Run Wizard of Oz sessions: You type responses pretending to be the AI.
  \item Test different ``personalities''---formal vs.\ friendly.
  \item Gather feedback on what users expect from the interaction.
\end{enumerate}

\textbf{What you learn}: Do users want step-by-step guidance? Do they get frustrated? Do they trust the responses?

\textbf{Only then}: Build/fine-tune the actual model.
\end{exbox}

\subsubsection*{Key Questions When Prototyping}

\begin{enumerate}
  \item \textbf{What assumptions are you testing?}
  \begin{itemize}
    \item User expectations?
    \item Technology feasibility?
    \item Trust and explainability?
  \end{itemize}
  
  \item \textbf{What does this mean for how/why you prototype?}
  \begin{itemize}
    \item Testing an idea? $\to$ Low fidelity.
    \item Testing specific interaction? $\to$ Vertical prototype.
    \item Testing overall flow? $\to$ Horizontal prototype.
    \item AI not built yet? $\to$ Wizard of Oz.
  \end{itemize}
\end{enumerate}

\begin{tipbox}[The Core Principle]
\textbf{Maximize the rate of learning by minimizing the time to try new ideas.}

Choose the \textit{minimum fidelity} needed to test your current assumptions. Don't over-engineer prototypes.
\end{tipbox}

\subsubsection*{Prototyping in the Design Process}

Recall the Double Diamond (Week 4):
\begin{center}
Discover $\to$ Define $\to$ \textbf{Develop} $\to$ Deliver
\end{center}

Prototyping sits in the \textbf{Develop} phase:
\begin{itemize}
  \item Diverge: Generate many prototype ideas.
  \item Converge: Test, iterate, refine toward final design.
\end{itemize}

In Stanford d.school terms: \textbf{Prototype} $\to$ \textbf{Test} (next week).

\subsubsection*{Worked Examples}

\begin{exbox}[Choosing Fidelity]
\textbf{Scenario}: You want to test whether users understand a new AI feature that explains why a recommendation was made.

\textbf{Early stage}: Paper prototype with explanation cards. Show to 5 users. Do they understand?

\textbf{Later stage}: Figma prototype with different explanation styles. A/B test which is clearer.

\textbf{Final stage}: Implement in code; run usability study.
\end{exbox}

\begin{exbox}[Vertical vs.\ Horizontal]
\textbf{Scenario}: AI-powered email assistant.

\textbf{Vertical prototype}: The ``compose reply'' feature works end-to-end (AI generates draft, user edits, sends). Other features (scheduling, search) are stubs.

\textbf{Horizontal prototype}: All features are visible and clickable, but none actually process data. Good for testing navigation: ``Can users find the AI features?''
\end{exbox}

\subsubsection*{Key Definitions Summary}

\begin{center}
\renewcommand{\arraystretch}{1.3}
\begin{tabular}{@{} l p{8cm} @{}}
\toprule
\textbf{Term} & \textbf{Definition} \\
\midrule
Prototype & Manifestation of design for stakeholder interaction; intentionally limited \\
Low-fidelity & Quick, cheap, conceptual; sketches, wireframes, paper \\
High-fidelity & Interactive, resembles final product; Figma, coded prototypes \\
Wireframe & Structural overview of layout and flow \\
Paper prototype & Hand-drawn sketch tested with users \\
Vertical prototype & Few features fully implemented (deep) \\
Horizontal prototype & All UI visible but shallow (no real processing) \\
Wizard of Oz & Human simulates AI behavior behind the scenes \\
\bottomrule
\end{tabular}
\end{center}

\subsubsection*{Tools Reference}

\begin{center}
\renewcommand{\arraystretch}{1.3}
\begin{tabular}{@{} l l l @{}}
\toprule
\textbf{Tool} & \textbf{Fidelity} & \textbf{Best For} \\
\midrule
Paper/Pen & Low & Very early ideation \\
Miro & Low & Collaborative wireframing \\
Balsamiq & Low & Quick wireframes \\
Figma & Low--High & Wireframes to interactive prototypes \\
Adobe XD & High & Interactive prototypes \\
Arduino/Pi & Physical & Tangible/IoT prototypes \\
\bottomrule
\end{tabular}
\end{center}

\subsubsection*{Recommended Resources}
\begin{itemize}
  \item \textbf{Textbook}: Chapter 11 (Prototyping)
  \item \textbf{Video}: Tony Fadell---Designing the Nest Thermostat
  \item \textbf{Tool}: Figma tutorials for interactive prototyping
\end{itemize}


